\section{What the error estimates are for}\label{ch12:sec:estimates}
\subsection*{Predicting how many steps are enough}
Take $T(x)=(x+1)/2$ on $\RR$, a contraction with factor $q=1/2$, and start from $x_0=0$. Then $x_1=1/2$, so $D=d(x_1,x_0)=1/2$, and the a priori estimate \eqref{eq:apriori} becomes
\[
|x_n-1|\le\frac{(1/2)^n}{1-1/2}\cdot\frac12=2^{-n}.
\]
Suppose we want an error below $0.001=1/1000$. Since $2^9=512$ and $2^{10}=1024$, the bound $2^{-n}$ is still larger than $1/1000$ at $n=9$ but equals $1/1024$ at $n=10$. So ten steps are guaranteed to be enough. Here the estimate cannot be improved: by the formula $x_n=1-2^{-n}$, the error is exactly $2^{-n}$, so nine steps really are too few. Usually, though, the true error is smaller than this worst-case bound, so the estimate may ask for more steps than necessary.

\subsection*{Certifying an approximation found by any method}
To certify a point $z$, the residual estimate \eqref{eq:residual} multiplies its residual by $1/(1-q)$, which can be large; the residual by itself is not a bound on the error. For example, $T(x)=0.99x+0.01$ is a contraction of $\RR$ with factor $q=0.99$, since $|T(x)-T(y)|=0.99\,|x-y|$, and its fixed point is $1$. The point $z=0$ has the small residual $|0-T(0)|=0.01$, yet its error is $|0-1|=1$, a hundred times larger. The residual estimate gives the bound $0.01/(1-0.99)=1$, exactly the true error. So a small residual guarantees a small error only when $q$ is not too close to $1$.

\subsection*{Using the length of the latest step}
For an iterate $x_n$ with $n\ge1$, there is a third bound, which uses the step just taken:
\[
d(x_n,x_*)\le\frac{q}{1-q}\,d(x_n,x_{n-1}).
\]
To prove it, apply the residual estimate to $z=x_n$. Its residual is $d(x_n,Tx_n)=d(x_n,x_{n+1})$, and since $x_n=Tx_{n-1}$ and $x_{n+1}=Tx_n$, the contraction inequality gives
\[
d(x_n,x_{n+1})=d(Tx_{n-1},Tx_n)\le q\,d(x_{n-1},x_n)=q\,d(x_n,x_{n-1}).
\]
So $d(x_n,x_*)\le d(x_n,x_{n+1})/(1-q)\le q\,d(x_n,x_{n-1})/(1-q)$.

This is called an \emph{a posteriori estimate}, from the Latin for ``from what comes after,'' because it can be evaluated only after the step from $x_{n-1}$ to $x_n$ has been computed. It is never worse than the a priori estimate, because Step~1 of the proof gives $d(x_n,x_{n-1})\le q^{n-1}D$, so
\[
\frac{q}{1-q}\,d(x_n,x_{n-1})\le\frac{q\cdot q^{n-1}D}{1-q}=\frac{q^n}{1-q}\,D,
\]
which is the a priori bound. When the observed steps shrink faster than the worst case $q^{n-1}D$, the a posteriori estimate is much better.

\pauseq{For $T(x)=(x+1)/2$ and $x_0=0$, compute the a posteriori bound for $x_3$ and compare it with the true error $|x_3-1|$.}{The iterates are $x_2=3/4$ and $x_3=7/8$, so the latest step has length $1/8$. With $q=1/2$ we have $q/(1-q)=1$, so the bound is $|x_3-1|\le1/8$. The true error is $1-7/8=1/8$, so the bound is attained, as the a priori bound $2^{-3}$ is for this map.}
